\documentclass[11pt]{ltjsarticle}
\usepackage{luatexja-fontspec}
\usepackage{amsmath,amssymb}
\usepackage[margin=22mm]{geometry}
\usepackage{graphicx}
\usepackage{longtable,booktabs,array}
\usepackage{calc}
\usepackage{xcolor}
\usepackage{hyperref}
\hypersetup{colorlinks=true,linkcolor=blue,urlcolor=blue,citecolor=blue}
\makeatletter
\def\maxwidth{\ifdim\Gin@nat@width>\linewidth\linewidth\else\Gin@nat@width\fi}
\makeatother
\setkeys{Gin}{width=\maxwidth,keepaspectratio}
\setcounter{secnumdepth}{0}
\setlength{\parindent}{0pt}
\setlength{\parskip}{0.5em}
\providecommand{\tightlist}{\setlength{\itemsep}{0pt}\setlength{\parskip}{0pt}}
\providecommand{\pandocbounded}[1]{#1}
\providecommand{\real}[1]{#1}
\newcounter{none}
\begin{document}
\section{補遺：離散性は外枠でなく節から生じる}\label{ux88dcux907aux96e2ux6563ux6027ux306fux5916ux67a0ux3067ux306aux304fux7bc0ux304bux3089ux751fux3058ux308b}

\subsection{------二値が生む節としての離散}\label{ux4e8cux5024ux304cux751fux3080ux7bc0ux3068ux3057ux3066ux306eux96e2ux6563}

\textbf{著者：} 木原 範昭（ORCID: 0009-0004-6753-4020） \textbf{日付：}
2026年6月18日 \textbf{DOI（本版）：} 10.5281/zenodo.20742991 ／
\textbf{Concept DOI（全版）：} 10.5281/zenodo.20742990 ／
\textbf{Zenodo:} https://zenodo.org/records/20742991
\textbf{関連論文：}「高い対称性の極限におけるスケールの対数写像とスケール・量子化の分離について」§3（Concept
DOI: 10.5281/zenodo.20740841）

\begin{center}\rule{0.5\linewidth}{0.5pt}\end{center}

\subsection{趣旨}\label{ux8da3ux65e8}

本補遺は、関連論文
§3「離散性の導出」を、独立に展開する。親論文に新しい主張を加えるものではなく、誤読されやすい一点------\textbf{離散性は外枠（有界性）からでなく、公理2（二値）が生む節から内在的に生じる}こと------を明示することを目的とする。

「離散性は有界性から出る」という説明は、弦の固有振動（両端を固定した弦には整数個の半波しか乗らない）を暗黙に借りている。だが端（境界）は片側にしか隣を持たない特権的な点であり、無名性に反する（親論文
補題）。端が無名性で禁じられる本系では、弦の境界条件を借りることはできない。にもかかわらず「有界だから離散」と述べると、こっそり弦の境界条件を持ち込むことになる------これは循環である。

本補遺は、この循環を避け、離散性が公理2（二値）から内在的に出ることを示す。離散の根は、外枠（有界性）ではなく、二値が必然的に生む節である。

本補遺は厳密な証明ではなく、親論文 §3
の構造の明示である。物理量への同定は行わない。

\begin{center}\rule{0.5\linewidth}{0.5pt}\end{center}

\subsection{1. 公理2
は節を生む}\label{ux516cux74062-ux306fux7bc0ux3092ux751fux3080}

本系の公理2（無名性の限界＝2）は、系の中の区別を二値------「他者なし」と「他者あり」------に畳む（親論文
§0.5・§2）。すなわち、系の中で立つのは二値（+/−、在る/無い）である。

ある量が二値（+/−）を取るとき、+ の領域から −
の領域へ移るには、必ず両者の境目を通る。この境目------符号が + から −
へ（あるいは − から +）切り替わる点------を\textbf{節}と呼ぶ。

二値であること自体が、節を生む。+ の領域と −
の領域が存在する以上、それらの間に必ず境目（節）がある。これは外から境界条件を課して生じるのではなく、二値（公理2）が内在的に生むものである。連続な量であれば符号の切り替わりは滑らかに通過する一点に過ぎないが、二値（公理2）の下では、+
と − の切り替わりが節として刻まれる。

\begin{center}\rule{0.5\linewidth}{0.5pt}\end{center}

\subsection{2.
節は離散である}\label{ux7bc0ux306fux96e2ux6563ux3067ux3042ux308b}

節は点である------+ と −
の境目という、特定の位置の点である。点は飛び飛びの対象であり、連続体ではない。

したがって、節が存在することは、即、離散的な対象が存在することを意味する。節と節の間には
+ または −
の領域があり、節そのものは離散的に並ぶ。離散性は、節という飛び飛びの点が存在することから、直ちに従う。

これは外枠（系の大きさの有界性）を必要としない。節は二値（公理2）から内在的に生じ、その節が飛び飛びの点であることから、離散性が出る。系の大きさに上限を課す必要はない。

\begin{center}\rule{0.5\linewidth}{0.5pt}\end{center}

\subsection{3.
無名性は節を整数個に縛る}\label{ux7121ux540dux6027ux306fux7bc0ux3092ux6574ux6570ux500bux306bux7e1bux308b}

節は点（+/−の切り替わり）である。無名性により、節には特権がない------すべての節は対等である（親論文の無名性）。

対等な節が並ぶとき、節の数は整数でなければならない。節は「在る」か「在らない」かの二値（公理2）であり、半端な節（1.5
個の節）はあり得ない。節は一つ、二つ、と数えられる離散的な対象であり、その総数は整数である。

すなわち、無名性が節を整数個に縛る。節の数が整数であることは、節が対等（無名）であり、かつ二値（在る/在らない）であることの帰結である。

\begin{center}\rule{0.5\linewidth}{0.5pt}\end{center}

\subsection{4.
閉じた系は節を偶数個に縛る}\label{ux9589ux3058ux305fux7cfbux306fux7bc0ux3092ux5076ux6570ux500bux306bux7e1bux308b}

系は閉じている（端を持たず一周して戻る、親論文
補題）。閉じた系の上で二値（+/−）が交互に並ぶとき、一周して元の符号（+）に戻るには、+
→ − → + → \ldots{} と、符号の切り替わりが偶数回起こらねばならない。

したがって、閉じた系では節の数は偶数である。一周して符号が戻るという閉性の要請が、節の数を偶数に縛る。

節と節の間が半波長に対応するから、節が偶数個であることは、許される波長（モード）が整数で番号づけられること（親論文
§3 第二段、λ\_n=L/n）に対応する。

\begin{center}\rule{0.5\linewidth}{0.5pt}\end{center}

\subsection{5. 導出の構造}\label{ux5c0eux51faux306eux69cbux9020}

以上より、離散性は外枠（有界性）でなく、二値（公理2）が生む節から内在的に出る。

{\def\LTcaptype{none} % do not increment counter
\begin{longtable}[]{@{}lll@{}}
\toprule\noalign{}
段階 & 根拠 & 結果 \\
\midrule\noalign{}
\endhead
\bottomrule\noalign{}
\endlastfoot
節の発生 & 公理2（二値、+/−の切り替わり） & 節が生じる \\
離散性 & 節は点（飛び飛び） & 離散的な対象が存在 \\
整数性 & 無名性（節が対等・二値） & 節の数は整数 \\
偶数性 & 閉性（一周で符号が戻る） & 節の数は偶数 \\
\end{longtable}
}

離散性の根は、二値（公理2）が生む節である。外枠（系の大きさの上限）は離散性のために必要ない。閉じた系（親論文
補題、自己証明不能の前提）の上に、二値が節を整数個・偶数個に刻むことで、許される波長（モード）が整数で番号づけられる（図1）。

\textbf{図1（fig5\_nodes\_on\_closed\_loop.svg）：}
閉じた系の上の節。二値（+/−）が交互に並び、その切り替わり点（節）が生じる。節は飛び飛びの点だから離散、無名性により整数個、閉性（一周で
+
に戻る）により偶数個に縛られる。離散性は外枠でなく、二値が生む節から出る。

\textbf{循環の回避：}
「有界だから離散」は弦の境界条件を借りる循環であった。本補遺の導出は、弦を借りず、節を公理2（二値）から内在的に導く。離散性は、系の大きさの有界性ではなく、二値という内在構造の帰結である。

\begin{center}\rule{0.5\linewidth}{0.5pt}\end{center}

\subsection{6.
限界（読み解きの注意）}\label{ux9650ux754cux8aadux307fux89e3ux304dux306eux6ce8ux610f}

本補遺は親論文 §3
の構造の明示であり、新しい主張を加えない。以下を明記する。

\begin{enumerate}
\def\labelenumi{\arabic{enumi}.}
\item
  \textbf{閉じた系は自己証明不能の公理（補題）である。}
  「無名性が端を禁じる→閉じる」の筋はあるが（親論文
  補題）、より基本的な原理から完全に導いたものではない。本補遺の偶数性（§4）は、この閉性を前提とする。
\item
  \textbf{離散性は公理2 と無名性から出る。}
  本補遺は離散性を外枠（有界性）から切り離し、公理2（二値）が生む節に帰した。ただし公理2
  自体の導出は親論文・本補遺の射程外である。
\item
  \textbf{「節」は系の中の二値の切り替わり点を指す。}
  物理的な波動の節への同定は行わない。二値（+/−、在る/無い）の切り替わりという形式的構造に限る。
\item
  \textbf{有界性は離散性のためには不要だが、モード番号の整数性には別途関わる。}
  親論文 §3 第二段は、閉じた系（周長 L）の上の定常波が整数モード
  λ\_n=L/n
  で番号づけられることを述べる。本補遺の節の偶数性（§4）は、これと対応する。有界性（系の大きさ）そのものは離散性の根拠ではないが、閉じた系の周長
  L はモードの番号づけに関わる。
\item
  \textbf{物理的同定を行わない。}
  二値・節・モードの物理的名称は無名のラベルであり、本補遺は特定の物理量への同定を行わない。
\end{enumerate}

\begin{center}\rule{0.5\linewidth}{0.5pt}\end{center}

\subsection{参照}\label{ux53c2ux7167}

\begin{itemize}
\tightlist
\item
  Kihara,
  N.「高い対称性の極限におけるスケールの対数写像とスケール・量子化の分離について------最小公理系からの限定的検証報告（改訂版）」§3、§2
  補題（閉性は無名性の帰結）、公理2。Concept DOI:
  10.5281/zenodo.20740841。
\end{itemize}
\end{document}
